\appendix
\section{Existence of detecting characters}\hypertarget{p107.sec.A}{}
We prove \hyperlink{p107.result.2.2}{Lemma 2.2} directly within the asymptotic-spectrum framework \cite{r16,r4}. A \emph{state} is a normalized, additive, restriction-monotone map $S\to\mathbb R_{\ge0}$; it is a character if it is also multiplicative. Every state satisfies $0\le\lambda(s)\le\mathrm R(s)$ and $\lambda(s)\ge1$ for $s\ne0$.

Fix integers $d\ge2$ and $k\ge0$ with $k<d^\nu$. To obtain a character with $\lambda(T_d)\ge k$, we will construct states satisfying the stronger family of inequalities
\begin{equation}
\lambda(T_ds)\ge k\lambda(s)\qquad(s\in S).
\tag{A.1}\hypertarget{p107.eq.A.1}{}
\end{equation}
These inequalities are preserved by normalized tensor multiplication: for any nonzero $z$, replacing $\lambda(x)$ by $\lambda(zx)/\lambda(z)$ preserves every inequality by applying \hyperlink{p107.eq.A.1}{(A.1)} to $zs$. The next lemma uses this invariance to obtain multiplicativity in a nonempty compact convex family of states. We then prove nonemptiness by deriving a contradiction to $k<d^\nu$ from an empty family.

\subsection{Obtaining multiplicativity}\hypertarget{p107.sec.A.1}{}
\begin{result}{Lemma}{A.1}{}
Let $\mathcal K$ be a nonempty compact convex set of states in the topology of coordinatewise convergence. Suppose that, for every nonzero $z\in S$, the map
\[
(P_z\lambda)(x)=\frac{\lambda(zx)}{\lambda(z)}
\]
takes $\mathcal K$ into itself. Then $\mathcal K$ contains a multiplicative state.
\end{result}
\begin{amsthmproof}
Since $\lambda(z)\ge1$, each $P_z$ is continuous in the product topology. The space $\mathbb R^S$ is Hausdorff and locally convex, so the Schauder--Tychonoff fixed-point theorem \cite{r17} gives a fixed point $\mu$ of $P_z$. Set $c=\mu(z)$. The subset
\[
\mathcal K'=\{\lambda\in\mathcal K:\lambda(zx)=c\lambda(x)\text{ for every }x\in S\}
\]
is nonempty and compact. Because $c$ is fixed, its defining equations are linear in $\lambda$, so it is convex. Taking $x=1$ shows that $\lambda(z)=c$ throughout $\mathcal K'$, so all its members are fixed by $P_z$. Moreover, every $P_w$ with $w\ne0$ preserves it:
\[
(P_w\lambda)(zx)=\frac{\lambda(zwx)}{\lambda(w)}=c\frac{\lambda(wx)}{\lambda(w)}=c(P_w\lambda)(x).
\]
Repeating the construction inside $\mathcal K'$ retains all earlier equations, so induction gives a common fixed point in $\mathcal K$ for every finite collection of maps $P_z$. Their fixed sets are closed, and compactness then gives a simultaneous fixed point for all nonzero $z$. It satisfies $\lambda(zx)=\lambda(z)\lambda(x)$ for $z\ne0$. Additivity implies $\lambda(0)=0$, which handles multiplication by zero.
\end{amsthmproof}

\subsection{Constructing the detecting state space}\hypertarget{p107.sec.A.2}{}
\begin{amsthmproof}[Proof of \hyperlink{p107.result.2.2}{Lemma 2.2}]
For the fixed integers $d,k$ above, let $\mathcal K$ be the subset of $\prod_{s\in S}[0,\mathrm R(s)]$ defined by normalization, additivity, monotonicity, and \hyperlink{p107.eq.A.1}{(A.1)}. The defining conditions are closed affine equalities and inequalities in a compact product, so $\mathcal K$ is compact and convex. It remains to prove nonemptiness before applying \hyperlink{p107.result.A.1}{Lemma A.1}.

\subsubsection*{Obtaining a catalytic comparison}
Suppose $\mathcal K$ were empty. By compactness, a finite subsystem of its defining conditions is inconsistent. Let $E\subset S$ contain every coordinate appearing in this subsystem, together with $0$ and $1$. Adjoin $\lambda(0)=0$, $\lambda(1)=1$, and write each coordinate bound homogeneously as
\[
\lambda(s)\ge0,\qquad\mathrm R(s)\lambda(1)-\lambda(s)\ge0\qquad(s\in E).
\]
The resulting finite system remains inconsistent.

To separate the equalities from the inequalities, let $e_s$, $s\in E$, be the coordinate vectors of $\mathbb R^E$. Quotient by the span of the selected additive relation vectors and $e_0$, and write $u$ for the image of $e_1$. Let $C$ be the cone generated by the images of all inequality vectors. Before taking the quotient, these vectors have the forms
\[
e_b-e_a\ (b\ge a),\qquad e_{T_ds}-ke_s,\qquad e_s,\qquad\mathrm R(s)e_1-e_s.
\]
The finitely generated cone $C$ is closed. If $-u\notin C$, separation gives a linear functional $L$ nonnegative on $C$ with $L(u)>0$. Dividing by $L(u)$ would give a normalized solution of the finite subsystem, including its bounds. This contradicts inconsistency, so $-u\in C$.

All vectors and relations have integer coefficients. Membership $-u\in C$ is a feasible rational linear system, so there are rational nonnegative cone coefficients and rational coefficients for the relation vectors. Clearing denominators yields a positive integer $m$ and an identity with nonnegative integer coefficients on the inequality vectors.

Map this identity to the additive Grothendieck group of $S$ by sending $e_s$ to $[s]$. This group adjoins formal additive inverses to the commutative monoid $(S,+)$, so all additive relation vectors map to zero. The bound vectors map to the order differences $[s]-[0]$ and $[\mathrm R(s)]-[s]$. Direct sums combine these and the other order differences into $[y]-[x]$ with $y\ge x$, while distributivity combines the terms from \hyperlink{p107.eq.A.1}{(A.1)} into $[T_ds]-k[s]$ for one $s\in S$. We obtain
\[
-m[1]=[y]-[x]+[T_ds]-k[s],\qquad y\ge x.
\]
Equality in this group means equality after adding a common element of the monoid. Thus there is $c\in S$ with
\[
m+y+T_ds+c=x+ks+c.
\]
Put $D=y+c$. Using $y\ge x$ gives the catalytic comparison
\begin{equation}
D+ks\ge D+m+T_ds.
\tag{A.2}\hypertarget{p107.eq.A.2}{}
\end{equation}
Here $s\ne0$: otherwise $D\ge D+m$, contradicting the additivity and monotonicity of any flattening rank.

\subsubsection*{Reusing the catalyst}
Iterating \hyperlink{p107.eq.A.2}{(A.2)} reuses the same tensor $D$, so its additive cost stays fixed as the number of copies grows. For every integer $n\ge1$,
\[
D+nks\ge D+nm+nT_ds\ge nT_ds.
\]
Since $s\ge1$, we have $\mathrm R(D)s\ge\mathrm R(D)\ge D$. With $K=nk+\mathrm R(D)$, it follows that
\begin{equation}
Ks\ge nT_ds,\qquad K^js\ge n^jT_d^js\qquad(j\ge1).
\tag{A.3}\hypertarget{p107.eq.A.3}{}
\end{equation}
For the induction, multiply the $j$th comparison by $K$, then use $Ks\ge nT_ds$:
\[
K^{j+1}s\ge n^jT_d^j(Ks)\ge n^{j+1}T_d^{j+1}s.
\]
This keeps a single factor $s$ and requires no cancellation. Also $K>0$, since the first right-hand side in \hyperlink{p107.eq.A.3}{(A.3)} is nonzero.

To bound $K$ using $\nu$, we convert the $n^j$ scalar copies in \hyperlink{p107.eq.A.3}{(A.3)} into a second matrix-multiplication factor. Fix $n$ and $\delta>0$. By the definition of $\nu$, choose a fixed integer $u\ge2$ with $\mathrm R(T_u)\le u^{\nu+\delta}$, and set
\[
\ell_j=\left\lfloor\frac{j\log_u n}{\nu+\delta}\right\rfloor,\qquad g_j=u^{\ell_j}.
\]
Submultiplicativity gives
\[
\mathrm R(T_{g_j})\le\mathrm R(T_u)^{\ell_j}\le u^{(\nu+\delta)\ell_j}\le n^j.
\]
Thus the diagonal tensor $n^j$ restricts to $T_{g_j}$. Multiplying this restriction by $T_{d^j}s$ and using $s\ge1$, we obtain
\[
K^js\ge n^jT_{d^j}s\ge T_{d^jg_j}s\ge T_{d^jg_j}.
\]
A direct sum of $K^j$ copies of $s$ has rank at most $K^j\mathrm R(s)$. Therefore
\[
\mathrm R(s)K^j\ge\mathrm R(T_{d^jg_j})\ge(d^jg_j)^\nu.
\]
First let $j\to\infty$, keeping $n,\delta,u$ fixed. Since $g_j^{1/j}\to n^{1/(\nu+\delta)}$, taking $j$th roots gives
\[
nk+\mathrm R(D)\ge d^\nu n^{\nu/(\nu+\delta)}.
\]
The choice of $u$ has disappeared from the inequality, so letting $\delta\downarrow0$ at fixed $n$ gives $nk+\mathrm R(D)\ge nd^\nu$. Dividing by $n$ and then letting $n\to\infty$ gives $k\ge d^\nu$, since the catalyst $D$ is fixed. This contradiction proves that $\mathcal K$ is nonempty.

To apply \hyperlink{p107.result.A.1}{Lemma A.1}, it remains to check that each $P_z$ preserves $\mathcal K$. For any nonzero $z$, the map preserves normalization, additivity, and monotonicity. It also preserves \hyperlink{p107.eq.A.1}{(A.1)}, because
\[
(P_z\lambda)(T_ds)=\frac{\lambda(T_dzs)}{\lambda(z)}\ge k\frac{\lambda(zs)}{\lambda(z)}=k(P_z\lambda)(s).
\]
Its coordinate bounds follow from $zx\le\mathrm R(x)z$. \hyperlink{p107.result.A.1}{Lemma A.1} supplies a multiplicative state in $\mathcal K$, hence a character. Taking $s=1$ in \hyperlink{p107.eq.A.1}{(A.1)} gives $\lambda(T_d)\ge k$.
\end{amsthmproof}
